% Generated draft: requires mathematical and source review.
\documentclass[11pt]{article}
\usepackage[margin=25mm]{geometry}
\usepackage{fontspec}
\setmainfont{Latin Modern Roman}
\usepackage{amsmath,amssymb,mathtools,unicode-math}
\setmathfont{Latin Modern Math}
\usepackage{enumitem,needspace,xurl,hyperref}
\setlength{\parindent}{0pt}
\setlength{\parskip}{5pt}
\setlength{\emergencystretch}{3em}
\newcommand{\N}{\mathbb{N}}
\newcommand{\Z}{\mathbb{Z}}
\newcommand{\Q}{\mathbb{Q}}
\newcommand{\R}{\mathbb{R}}
\newcommand{\C}{\mathbb{C}}
\newcommand{\F}{\mathbb{F}}
\newcommand{\A}{\mathbb{A}}
\newcommand{\PP}{\mathbb{P}}
\newcommand{\E}{\mathbb{E}}
\newcommand{\eps}{\varepsilon}
\newcommand{\dd}{\,\mathrm{d}}
\begin{document}
\section*{UnsolvedMath ID: 800692}
\textbf{Source problem code:} catalogue entry
\subsection{Definitions and mathematical statement}
Let \(p\) be a prime and let \(F\) be a field of characteristic \(p\). A central simple \(F\)-algebra is a finite-dimensional associative unital \(F\)-algebra whose center is \(F\) and which has no nonzero proper two-sided ideals. Its degree is \(\deg(A)=\sqrt{\dim_F A}\). Under the usual characteristic-\(p\) convention, a \(p\)-algebra is a central simple algebra of degree a power of \(p\).

For a central division \(F\)-algebra \(D\), being a crossed product means that \(D\) contains a maximal commutative subfield \(L\) such that \(L/F\) is Galois; necessarily \([L:F]=\deg(D)\).

\textbf{Question, under the division-algebra interpretation.} Is it true that for every prime \(p\), every field \(F\) of characteristic \(p\), and every central division \(F\)-algebra \(D\) of degree \(p^2\), there exists a subfield \(L\subseteq D\) containing \(F\) such that
\[
[L:F]=p^2\qquad\text{and}\qquad L/F\text{ is Galois}?
\]
No cyclicity condition is imposed on \(\operatorname{Gal}(L/F)\).

\textbf{Unresolved scope qualification.} The supplied question says only ``\(p\)-algebras of degree \(p^2\)'' and does not reproduce the preceding definitions. In particular, it does not explicitly restrict to division algebras or specify the crossed-product convention for general central simple algebras. The displayed question gives the division-algebra interpretation; confirmation of the source's conventions is required before identifying it unconditionally with the intended statement.

The supplied background reports that the smallest known non-crossed-product \(p\)-algebras have degree \(p^3\), and presents the degree-\(p^2\) question as open. These are supplied historical source qualifications, not independently verified assertions about current knowledge.
\subsection{Short English statement}
Under the division-algebra interpretation, the question asks whether every central division algebra of degree \(p^2\) over every field of characteristic \(p\) contains a Galois maximal subfield. The supplied excerpt leaves the precise algebra and crossed-product conventions unresolved.
\subsection{Sources}
\begin{itemize}
\item Local UnsolvedMath registry, permanent ID \(800692\), exact source problem code \texttt{catalogue entry}, title \emph{Open Problems on Central Simple Algebras --- Problem 2}. The supplied registry labels the problem open and attributes it to Auel, Brussel, Garibaldi, and Vishne, with proposed year 2010; this status is supplied metadata, not an independently verified mathematical assessment.
\item Auel, Brussel, Garibaldi, and Vishne, \emph{Open Problems on Central Simple Algebras} (2010), Problem 2. Supplied source URL: \url{https://arxiv.org/abs/1006.3304}. The record reports access on 2026-10-16 and extraction by \texttt{arxiv\_latex}; these are supplied provenance details, not a claim of independent source access or verification.
\end{itemize}
\end{document}
